We present MAgHARCM (Multi-Agent Harness for Archaeology-guided Repository-level Code Migration), a Go implementation of the ReCodeAgent architecture that runs entirely against locally-hosted small language models (SLMs) through Ollama. MAgHARCM instantiates the four-agent Analyzer / Planning / Translator / Validator cyclic graph from ReCodeAgent on top of CloudWeGo Eino, adds a coverage-guided characterization test pass, and exposes a pluggable Language Server Protocol (LSP) provider that can fall back to native tree-sitter parsing when ABCoder MCP is unavailable. We exercise the harness on the same public corpora used by ReCodeAgent, AlphaTrans, CodePlan, and Oxidizer (C, Go, and Rust source projects drawn from GildedRose, Oxidizer, and the AlphaTrans subject set) and report observed compile and test outcomes directly from \texttt{.artifacts/} run logs. Sample 1 (GildedRose C$\rightarrow$Rust) reaches 100\% compilation across all ten repair iterations but 3 of 7 translated tests fail on edge-case inventory rules (negative \texttt{sell\_in} backstage pass, Sulfuras quality preservation, and elapsed-time backstage pass) -- a 57\% test pass rate that we attribute to a hallucinated \texttt{quality = 50} branch rather than the source's zero-out branch. Sample 2 (Oxidizer stats Go$\rightarrow$Rust) saturates the harness: the single-shot Translator produces zero files because the Planning agent falls back to a two-line skeleton for 517 fragments, which is the failure mode that motivates the chunked Translator primitive (NEW-PRIM-23) we describe. Samples 3 and 4 are deferred until that primitive lands. The paper documents both the successes and the failure modes as evidence that the four-agent decomposition carries over to a local-SLM backend, but that the Translator step needs fragmentation for repositories whose fragment count exceeds what fits in a single coding-model prompt.
